% This is samplepaper.tex, a sample chapter demonstrating the
% LLNCS macro package for Springer Computer Science proceedings;
% Version 2.21 of 2022/01/12
%
\documentclass[runningheads]{llncs}
%
\usepackage[T1]{fontenc}
% T1 fonts will be used to generate the final print and online PDFs,
% so please use T1 fonts in your manuscript whenever possible.
% Other font encondings may result in incorrect characters.
%
\usepackage{graphicx}
\usepackage{todonotes}
\usepackage{amsmath,amssymb}
% Used for displaying a sample figure. If possible, figure files should
% be included in EPS format.
%
% If you use the hyperref package, please uncomment the following two lines
% to display URLs in blue roman font according to Springer's eBook style:
%\usepackage{color}
%\renewcommand\UrlFont{\color{blue}\rmfamily}
%\urlstyle{rm}
%
\begin{document}
%
\title{Volumetric Gaussian Splatting: Direct Optimization in Voxel Space for Efficient Medical Data Visualization}
%
\titlerunning{Volumetric Gaussian Splatting}
% If the paper title is too long for the running head, you can set
% an abbreviated paper title here
%
\author{Valentin Kraft\inst{1}\orcidID{0000-0003-4057-2452} \and
Andrea Schenk\inst{1}\orcidID{0000-0003-1806-8380} \and
Gabriel Zachmann\inst{2}\orcidID{0000-0001-8155-1127}}
%
\authorrunning{V. Kraft et al.}
% First names are abbreviated in the running head.
% If there are more than two authors, 'et al.' is used.
%
\institute{Fraunhofer MEVIS, 28359 Bremen, Germany \and
Computer Graphics and Virtual Reality Department, University of Bremen, 28359 Bremen, Germany\\}
%
\maketitle              % typeset the header of the contribution
%
\begin{abstract}
We present a novel formulation of 3D Gaussian Splatting for volumetric data that departs from the conventional paradigm of multi-view image supervision. Instead of optimizing against rendered RGB images, our approach directly operates in voxel space using target volumes and region-of-interest constraints, enabling a volume-native optimization process. To this end, we introduce a differentiable splat-to-volume rasterization that projects anisotropic Gaussians into a voxel grid and allows direct optimization of spatial and radiometric parameters.

Each Gaussian explicitly encodes a physical intensity value, such as Hounsfield units in computed tomography, enabling a volume-rendering-like representation with post hoc transfer function control. This allows flexible visualization without retraining and establishes a closer connection to classical volume rendering pipelines.

Our method further incorporates volume-driven initialization strategies and stable optimization mechanisms tailored to volumetric supervision. We demonstrate that it enables significantly more efficient rendering, making it particularly suitable for resource-constrained environments such as mixed reality devices.

Overall, our work reframes Gaussian Splatting as a direct volumetric representation and optimization problem, opening new directions for efficient and controllable visualization of medical data.

\keywords{Gaussian Splatting \and Medical Visualization \and Volume Rendering \and Differentiable Rendering \and Volumetric Optimization}
\end{abstract}
%
%
%

\section{Introduction}

Recent advances in medical image analysis have led to highly accurate and increasingly automated segmentation of anatomical structures, driven by methods such as TotalSegmentator \cite{totalsegmentator} and foundation models like SAM \cite{sam}. While these developments enable rich semantic understanding of volumetric data, efficient and high-quality visualization of such data—particularly in immersive environments—remains a significant challenge.

Volume rendering is the de facto standard for visualizing medical data, as it preserves the full volumetric information and allows flexible exploration via transfer functions. However, achieving high visual fidelity typically requires substantial computational resources, which limits its applicability in resource-constrained settings such as mixed reality (MR) devices. Existing solutions often rely on streaming-based approaches, where rendered images or intermediate representations are transmitted from powerful servers to lightweight clients. While effective, these approaches introduce system complexity, latency, and potential failure modes that are undesirable in clinical or interactive scenarios.

At the same time, 3D Gaussian Splatting \cite{kerbl2023gaussiansplatting} has recently emerged as a highly efficient rendering paradigm, enabling real-time visualization of complex scenes with impressive performance characteristics. However, existing formulations are fundamentally designed for image-based reconstruction, where a scene is learned from multi-view RGB observations under known or estimated camera parameters. This paradigm introduces a complex processing pipeline—including camera pose estimation, feature extraction, and multi-view consistency constraints—and defines supervision in the image domain.

While this formulation is well suited for natural image data, it is not well aligned with volumetric domains such as medical imaging. In these settings, ground-truth data is directly available in voxel space, providing a consistent spatial reference frame and physically meaningful intensity values (e.g., Hounsfield units). Optimizing a representation through rendered images therefore introduces an unnecessary level of indirection and reliance on camera-based supervision introduces additional sources of error and complexity that are unneccesary in the volumetric setting.

These limitations motivate a reformulation of Gaussian Splatting that operates directly in the native domain of the data. In this work, we revisit 3D Gaussian Splatting from a volume-centric perspective and cast it as a direct volumetric optimization problem. Instead of relying on image-based supervision, we optimize a set of anisotropic Gaussians against a target volume, guided by region-of-interest constraints. This eliminates the need for intermediate rendering during training and removes dependencies on camera models, feature extraction, and pose estimation.

A key aspect of our approach is an intensity-aware Gaussian representation, where each splat encodes a physical intensity value. This allows the learned representation to behave similarly to classical volume data, enabling flexible post hoc transfer function design without retraining. To support this formulation, we introduce a differentiable splat-to-volume rasterization that enables end-to-end optimization of Gaussian parameters directly in voxel space, along with volume-driven initialization and stabilization strategies tailored to this setting.

While our method does not aim to fully replace high-quality volume rendering in terms of visual fidelity, it offers a compelling trade-off between quality and efficiency. In particular, it enables fast, lightweight rendering that is well suited for interactive and immersive applications where computational resources are limited.

%In summary, our work bridges the gap between Gaussian Splatting and volumetric medical data by reformulating the method as a direct optimization problem in voxel space. This results in a representation that is more aligned with the underlying data domain, simplifies the reconstruction pipeline, and enables efficient and flexible visualization of volumetric data. \textbf{TODO: LÖSCHEN}?

Our main contributions are as follows:

\begin{itemize}

\item \textbf{Volume-supervised 3D Gaussian Splatting.}
We reformulate 3D Gaussian Splatting as a volume-native optimization problem by replacing conventional multi-view RGB supervision with direct supervision in voxel space using target volumes and region-of-interest masks. This is enabled by a differentiable splat-to-volume rasterization that allows end-to-end optimization of Gaussian parameters in the volume domain.

\item \textbf{Intensity-aware volumetric Gaussian representation.}
We introduce a volume-native Gaussian representation in which each splat encodes a physical intensity value (e.g., Hounsfield units). This allows the model to behave similarly to classical volume data, enabling post hoc modification of transfer functions without retraining.

\item \textbf{Efficient real-time volumetric rendering.}
We demonstrate that the proposed representation provides a compact and efficient alternative to dense voxel grids, enabling real-time, volume-rendering-like visualization suitable for resource-constrained environments such as mixed reality devices.

\end{itemize}


\section{Related Work}

\paragraph{Gaussian Splatting and Neural Rendering:}
Recent advances in neural rendering have introduced highly efficient scene representations, among which 3D Gaussian Splatting \cite{kerbl2023gaussiansplatting} has emerged as a particularly impactful approach, sparking a lot of follow-up research. The foundational work represents scenes as sets of anisotropic Gaussians optimized from multi-view image supervision, enabling extremely efficient real-time rendering with high visual fidelity.

Subsequent work has explored multiple extensions of this paradigm. Siemens Healthineers proposed "Cinematic Gaussians" \cite{cinematicgaussians2024}, which demonstrates that Gaussian Splatting can achieve near photorealistic visualization of anatomical structures using advanced rendering techniques such as path tracing. Similarly, Zhang et al. \cite{relightablegaussian2024} introduces relightable representations that further reduce the need for traditional volume rendering.

Recent work also investigates more advanced formulations of Gaussian-based rendering. Gao et al. \cite{ddgsct2024,sixdgs2024} extend Gaussian Splatting toward volumetric rendering scenarios by incorporating direction-dependent effects and improving fine-detail reconstruction. Also interaction techniques which are known from volume rendering, such as clipping \cite{raraclipper2025} or transfer-function-like behaviour \cite{ivrgs2025} are beginning to get transferred to Gaussian Splatting.

Despite these advances, all existing approaches remain fundamentally tied to image-based supervision and view-dependent rendering. Even methods targeting volumetric effects still rely on rendering-based optimization rather than directly leveraging volumetric ground-truth data.



\paragraph{Volume Rendering and Splatting Techniques:}

Volume rendering is the standard approach for visualizing volumetric data and has been extensively studied in computer graphics. Early work \cite{drebin1988volume} established the foundations for direct volume visualization, while subsequent approaches have focused on improving realism and interactivity. For instance, Path Tracing \cite{kroes2012exposure} demonstrates high-quality photorealistic rendering using physically-based techniques.

Early splatting-based approaches, such as "EWA Volume Splatting" \cite{ewasplatting2001}, approximate volume rendering using Gaussian kernels and share conceptual similarities with modern Gaussian Splatting methods. More recent discussions \cite{cgf2025gs} highlight ongoing challenges in achieving both high fidelity and efficiency.

Despite significant progress, classical and modern volume rendering approaches remain computationally demanding, particularly for large scale medical datasets and real time applications.



\paragraph{Efficient Medical Visualization and Mixed Reality:}


Efficient visualization of volumetric medical data is critical for interactive and immersive applications. Surveys \cite{mrreview2024,mrreview2} highlight both the potential and limitations of mixed reality in clinical workflows.

Prior work has explored optimized rendering techniques \cite{avis2024}, which demonstrate that improved volume rendering can positively impact clinical decision-making. However, even optimized approaches may struggle to meet performance requirements on resource-constrained devices without additional system complexity.


% \subsection{Segmentation and Structure-Aware Medical Imaging}


% Recent advances in medical image analysis have enabled accurate and large-scale segmentation of anatomical structures. Methods such as nnU-Net \cite{nnunet}, TotalSegmentator \cite{totalsegmentator}, and Segment Anything Model \cite{sam} provide detailed structural information directly from volumetric data.

%However, most visualization pipelines do not integrate this structural information at the representation level, limiting efficiency and flexibility.


\paragraph{Summary:}

Gaussian Splatting has enabled highly efficient rendering but remains fundamentally image-based, even in recent volumetric extensions. At the same time, classical volume rendering provides accurate visualization at high computational cost. Medical imaging, however, offers direct access to volumetric data and structural annotations. Our work bridges this gap by reformulating Gaussian Splatting as a direct volumetric optimization problem, enabling efficient and flexible visualization without intermediate rendering.




\section{Method}

\subsection{From Image-Space to Volume-Space}
We reformulate 3D Gaussian Splatting as a direct volumetric optimization problem. Instead of learning a scene representation from multi-view RGB images, we optimize a set of anisotropic 3D Gaussians directly against a target volume $V$ and a corresponding region-of-interest (ROI) mask $M$. As a result, supervision, initialization, and topology adaptation are all performed in voxel space, transforming Gaussian Splatting into a representation that is naturally aligned with medical imaging and other volumetric domains. This reformulation also significantly reduces the complexity of the pipeline and eliminates the need for camera models, feature extraction, feature matching and image-based supervision.

Given $V$ and $M$, the objective is to learn a compact Gaussian set that approximates the underlying volumetric signal within relevant regions. Optimization is performed entirely in voxel space via a differentiable splatting formulation.


\begin{figure}
\centering
\includegraphics[width=0.9\textwidth]{viewer_app.png}
\caption{Our novel gs viewer exploits the new per-splat intensity values (effectively storing HU values) to dynamically change a transfer function during rendering, similar to a volume rendering.} \label{fig:app}
\end{figure}


\subsection{Gaussian Representation}
Similar to the original 3DGS approach, we parameterize each Gaussian $i$ by its center position $\mu_i$, its covariance matrix $\Sigma_i$ and opacity $\alpha_i$. In addition to that, each splat explicitly encodes a physical intensity value $c_i$ (e.g., Hounsfield units in CT), sampled directly from the input volume at each Gaussian position using trilinear interpolation in normalized volume space. For large splats, the method can instead use the mean of the voxels covered by the splat footprint.

This allows transfer functions to be applied or modified after training (see Fig \ref{fig:app}), similar to classical volume rendering, without requiring re-optimization. In contrast to the original method, however, we use a spherical harmonic degree of 0 or 1 to avoid view-dependent colors, which would not fit to the expected precision in medical applications.

In order to still be somewhat compatible with other gaussian splatting tools and workflows, the exported PLY files contain a grayscale sh color value that gets computed by converting the intensity values to a SH color using a standard black-and-white transfer function.


\subsection{Differentiable Splat-to-Volume Formulation}
To compare the Gaussian representation to the target data, we rasterize the current Gaussian set into a dense voxel grid defined over the volume or ROI bounds. For a voxel center $\mathbf{x}$, Gaussian $i$ contributes an opacity-weighted anisotropic kernel in its local frame,

\begin{equation}
w_i(\mathbf{x}) = \alpha_i \exp\left(-\tfrac{1}{2}\left\|S_i^{-1} R_i^\top (\mathbf{x}-\mu_i)\right\|_2^2\right),
\end{equation}

where $\mu_i$ is the Gaussian center, $R_i$ the rotation, $S_i$ the diagonal scale matrix, and $\alpha_i$ the opacity. For intensity supervision, the voxel value is obtained as a normalized weighted average of the per-Gaussian intensities $c_i$,

\begin{equation}
\hat{V}(\mathbf{x}) = \frac{\sum_i w_i(\mathbf{x}) c_i}{\sum_i w_i(\mathbf{x}) + \varepsilon},
\qquad
\hat{D}(\mathbf{x}) = 1 - \exp\left(-\gamma \sum_i w_i(\mathbf{x})\right)
\end{equation}

for intensity and density rendering, respectively. In practice, a dense all-to-all evaluation would be too expensive, so our implementation uses several bounded approximations: Gaussians are processed in batches, each splat is evaluated only inside a finite local support derived from a kernel cutoff, the support radius is capped in voxel units, very small splats are clamped to a minimum voxel-space standard deviation, and for large point sets accumulation can be performed on a downscaled working grid that is trilinearly upsampled to the target resolution. This sparse local accumulation keeps the rasterization differentiable with respect to positions, scales, rotations, opacities, and intensities while making training feasible for larger volumes.

\begin{quote}
\noindent\textbf{High-level pseudo code:}
\begin{tabbing}
\hspace{1.2em}\=\hspace{1.2em}\=\kill
initialize numerator grid $A \leftarrow 0$ and weight grid $W \leftarrow 0$ \\
optionally choose a downscaled working grid \\
\textbf{for} each Gaussian batch $\mathcal{B}$ \textbf{do} \\
\> map Gaussian centers to working-grid coordinates \\
\> determine a finite voxel neighborhood for each Gaussian \\
\> evaluate anisotropic kernel weights inside that neighborhood \\
\> accumulate $A \mathrel{+}= w_i c_i$ and $W \mathrel{+}= w_i$ \\
\textbf{end for} \\
\textbf{if} intensity rendering \textbf{then} output $\hat{V} = A / (W + \varepsilon)$ \\
\textbf{else} output $\hat{D} = 1 - \exp(-\gamma W)$ \\
\textbf{if} a working grid was used \textbf{then} trilinearly upsample to full resolution
\end{tabbing}
\end{quote}

\subsection{Initialization}

For almost all clinical tasks, e.g. for intervention planning, only specific structures are relevant. Therefore, and to focus model capacity and compute power on only where it is necessary, we assume the availability of a ROI mask. The ROI mask could for instance come from a machine learning-based segmentation model like TotalSegmentator and can be a binary or soft (float) mask and could cover one or multiple relevant structures. During initialization, the Gaussian spawn-positions are sampled inside the ROI mask. For soft masks, sampling probabilities correspond to the voxel value (which, depending on the segmentation method, might be interpreted as a segmentation confidence), allowing the method to retain uncertain but potentially important regions.

In addition, we incorporate simple structural priors derived from the input volume. Local volume gradients and second-order cues (e.g., structure tensor or Hessian-based estimates) are used to estimate preferred orientations, enabling anisotropic Gaussians to align with elongated structures such as vessels. This provides a strong initialization that improves convergence and reduces the need for extensive topology adaptation during training.

The initial splat scale is sampled per Gaussian from a uniform range defined in voxel units, using the mean voxel spacing as the isotropic reference scale:

\begin{equation}
s_i=\bar v\left(\sigma_{\min}^{\mathrm{vox}}+u_i\left(\sigma_{\max}^{\mathrm{vox}}-\sigma_{\min}^{\mathrm{vox}}\right)\right).
\end{equation}

Here, $s_i$ is the initial base scale of Gaussian $i$, $\bar v$ is the mean voxel size, $\sigma_{\min}^{\mathrm{vox}}$ and $\sigma_{\max}^{\mathrm{vox}}$ are the user-defined minimum and maximum initialization scales in voxel units, and $u_i \sim \mathcal{U}(0,1)$ is a random sample that selects a scale uniformly within that interval.

\subsection{ROI-Aware Optimization}
Training is performed directly in voxel space using losses restricted to the ROI. This prevents large empty background regions from dominating optimization and improves robustness for sparse anatomical structures.

The framework supports both structure-focused optimization (e.g., mask fitting) and intensity-based optimization, as well as their combination. By default, the loss for the structure-focused optimization is Dice and the standard loss for the intensity / CT data values is mean squared error (MSE). An additional penalty discourages density leakage outside the ROI, which stabilizes training.

\begin{equation}
\mathcal{L} = \lambda_{\mathrm{mask}} \mathcal{L}_{\mathrm{mask}} + \lambda_{\mathrm{ct}} \mathcal{L}_{\mathrm{ct}} + \lambda_{\mathrm{leak}} \mathcal{L}_{\mathrm{leak}}.
\end{equation}

\subsection{Adaptive Topology Optimization}
To maintain an efficient representation, the Gaussian set is dynamically adapted during training. Densification and pruning are performed periodically during a predefined training window. At each densification step, Gaussians with large accumulated position gradients (measured via the norm of voxel-space position updates) are refined: small splats are cloned, larger splats are split, and low-opacity splats are removed. The selection threshold is set adaptively from the current gradient distribution, and an additional hole-filling step can spawn new splats in under-covered regions (e.g., regions with low accumulated weights in the voxel grid). Unlike standard Gaussian Splatting, these topology updates are driven entirely by voxel-space optimization signals rather than by view-dependent visibility.

This enables effective modeling of thin and complex anatomical structures while keeping the representation compact.

\subsection{Implementation Details}
Our implementation follows the original 3D Gaussian Splatting pipeline and is based on the original implementation, but replaces image-space supervision with a dedicated volumetric training path. In particular, data loading, initialization, supervision, rasterization, optimization, and export are integrated into a single training workflow that operates directly on target volumes and ROI masks. Data can be loaded using the NIfTI file format.

To make training practical for large medical volumes, the implementation includes several memory- and runtime-aware design choices, including ROI-restricted computation, optional downscaling of the working grid, chunked accumulation of Gaussian contributions, and mixed-precision storage where appropriate. We further use monitoring and logging utilities to track optimization stability and support debugging during long training runs.

The system also includes task-specific presets for different anatomical regimes (e.g. for vessels) and exports intermediate checkpoints and PLY snapshots for inspection and downstream rendering. Overall, the implementation is designed to keep volumetric optimization tractable while preserving the flexibility needed for different datasets, supervision modes, and hardware constraints. All hyperparameters are made accessible as CLI parameters. For optimization, we employ the Adam optimizer.

Since standard Gaussian splatting viewers do not support our new intensity value, which we store per splat, we created a novel viewer application that allows dynamic changes of the transfer function during rendering, similar to a volume rendering workflow (see Fig \ref{fig:app}). This functionality is essential to fully exploit the intensity-aware representation.

The implementation (including the novel viewer application) is publicly available at \url{https://github.com/your-repo/3d-gaussian-splatting}


% ----------------------------

\section{Experiments}

\subsection{Experimental Setup}
We evaluate our method on medical volumetric datasets, focusing on anatomically relevant structures such as vessels and organs. The input data consists of CT images together with corresponding segmentation masks.

We use the ``Case 1'' dataset from the AbdomenCT-1k project \cite{ma2021abdomenct}. The dataset was resampled to an isotropic voxel spacing of 1.5\,mm and size $292 \times 292 \times 180$. Segmentations were obtained using TotalSegmentator v2 and saved as binary masks. In addition, a soft liver segmentation mask was created using an internal segmentation model.

All experiments are conducted on a Windows 11 system with an Intel i7 Ultra 155H, 32GB RAM, and an NVIDIA RTX 4060 Mobile GPU (8GB VRAM).

\subsection{Baselines}
We compare our method against the following baselines:

\begin{itemize}
\item \textbf{Volume Rendering (VR):} Standard GPU-based volume rendering using the GigaVoxel renderer in MeVisLab 4.2.70 \cite{mevislab}. This serves as a reference for visual fidelity.

\item \textbf{3D Gaussian Splatting (3DGS):} Standard image-based 3DGS as implemented in Nerfstudio \cite{nerfstudio}. Training is performed on 500 rendered images with a resolution of $1024 \times 1024$ and from random spherical camera positions. The renderings were generated with the same GigaVoxel renderer to ensure consistency of input data.
\end{itemize}

\subsection{Evaluation Protocol}
\label{sec:evaluation}

Our evaluation protocol is designed to reflect the inherent differences between our method and standard 3D Gaussian Splatting (3DGS). While both approaches rely on a Gaussian scene representation, they are trained under fundamentally different supervision regimes. Our method is optimized directly in the volume domain using a target volume and ROI mask, whereas standard 3DGS is trained from multi-view images under known or estimated camera parameters. Consequently, the methods differ in input modality, optimization objective, coordinate system, and computational structure, and thus address related but inherently different reconstruction problems.

This discrepancy is further accentuated by our active-set training strategy, where only a subset of Gaussians is updated per iteration (controlled via \texttt{max\_points\_per\_iter}), while evaluation is always performed on the full set. As a result, direct quantitative comparisons between both approaches are not well-defined and may lead to misleading conclusions. In particular, quantitative comparisons in image space would be confounded by differences in supervision, rendering assumptions, and pose alignment.

We therefore adopt a two-part evaluation protocol consisting of (i) a quantitative evaluation in the volume domain and (ii) a qualitative comparison in image space.

\paragraph{Quantitative evaluation:}
Our primary quantitative metric is the masked mean squared error (MMSE) computed in the native volume domain. After training, the Gaussian representation is rasterized onto the original voxel grid of the input volume. The reconstructed volume is then compared to the ground-truth volume within the region of interest (ROI), defined by a binary mask. This ensures that evaluation is performed in the same domain as supervision, avoiding ambiguities related to camera placement or cross-method alignment. Importantly, all Gaussians are used during evaluation, irrespective of the subset-based updates during training.

Formally, let $V$ denote the normalized ground-truth volume, $\hat{V}$ the reconstructed volume obtained by rasterizing the Gaussian model, and $M \in \{0,1\}$ the ROI mask. The MMSE is defined as

\begin{equation}
\mathcal{L}_{\mathrm{MMSE}} =
\frac{1}{\sum_i M_i}
\sum_i M_i \, (\hat{V}_i - V_i)^2,
\end{equation}

where the sum runs over all voxels and only voxels within the ROI contribute to the error. This metric serves as our primary quantitative measure of reconstruction fidelity.

\paragraph{Qualitative evaluation:}
In addition to the volume-domain evaluation, we perform a qualitative comparison against direct volume rendering of the input data and against standard 3DGS (see Figure \ref{fig:comparison}). Since these methods operate in different coordinate systems and are not naturally aligned to a shared camera model, renderings are generated from approximately matched viewpoints rather than enforcing exact pose correspondence. These comparisons are intended solely for visual inspection of structural fidelity, anatomical plausibility, and perceptual quality, and are not used for quantitative evaluation.


\section{Results}



\subsection{Quantitative Results}
We report the MMSE values for our method as defined in Sec.~\ref{sec:evaluation} in Table \ref{tab:mmse}. Please note that MMSE values for standard 3DGS cannot be computed in a meaningful way without an additional registration step or similar calibration steps, as the resulting models produced by its image-space optimization are usually not aligned with the original voxel grid / coordinate system.

\begin{table}
\centering
\caption{MMSE and ROI-based PSNR for our volume-supervised method for iteration 1000, evaluated against the ground-truth input volume. The Gaussian model is rasterized onto the native evaluation grid, and the metrics are computed only within the ROI mask.}
\label{tab:mmse}
\begin{tabular}{|l|c|c|}
\hline
Splat Count & MMSE & PSNR [dB] \\
\hline
200k & 0.002680 & 25.72 \\
500k & 0.002458 & 26.09 \\
1M &   0.002126 & 26.72 \\

\hline
\end{tabular}
\end{table}

The low MMSE values obtained in our experiments indicate that the proposed method reconstructs the target volume with high fidelity inside the ROI. Although this metric provides a precise volume-domain measure, it is complemented by qualitative inspection to assess fine-scale structural accuracy beyond the averaged error. \todo{In discussion!}

\subsection{Qualitative Results}

\begin{figure}
\includegraphics[width=\textwidth]{comparison.png}
\caption{The abdomen dataset rendered using different modalities. Top left: Volume Rendering (reference), top right: Standard 3DGS (253k Splats), bottom left: our method (200k Splats), bottom right: our method (1M Splats). It is obvious that the standard 3DGS model appears a bit sharper, but less precise with respect to the colors.} \label{fig:comparison}
\end{figure}

We qualitatively compare our method to direct volume rendering (reference) and to standard 3DGS. Figure \ref{fig:comparison} shows representative renderings from approximately matched viewpoints.

Our method preserves the overall anatomical structure while significantly reducing representational complexity. In particular, fine structures such as vessels remain visually consistent, whereas 3DGS exhibits artifacts due to its reliance on image-based supervision.\todo{In discussion!}

\subsection{Efficiency Analysis}

\paragraph{Rendering performance:}
We evaluate rendering performance in terms of frames per second (FPS). Table \ref{tab:fps} compares our method to volume rendering and 3DGS. We used the "GigaVoxelRenderer" in MeVislab for measuring the Volume Rendering. For 3DGS and our method, we used the "SIBR Gaussian Viewer App" that is part of the standard 3DGS repository. All measurements were taken at a rendering resolution of $1024 \times 1024$.

\begin{table}
\centering
\caption{Rendering performance of the different methods. Note that the rendering FPS is bounded by the SIBR viewer implementation and hardware constraints (240Hz display) and therefore does not fully reflect the theoretical performance differences between methods.}\label{tab:fps}
\begin{tabular}{|l|c|c|}
\hline
Method & Nr. of primitives & Dataset 1 (Abdomen)\\
\hline
Volume Rendering & 15.347.520 & 107 FPS\\
3DGS & 200k & 240 (max) FPS \\
Ours & 200k & 240 (max) FPS \\
Ours & 500k & 240 (max) FPS \\
Ours & 1M & 147 FPS \\
\hline
\end{tabular}
\end{table}

\paragraph{Training performance:}
We also present the timings for model training using our method in table \ref{tab:training}. Since our method significantly differs from standard Gaussian Splatting (see section \ref{sec:evaluation}), we do not compare against it in terms of training times.

Training time depends not only on the number of Gaussians but also on the chosen subset size, voxel resolution, and ROI extent. The reported values should therefore be interpreted as indicative rather than directly comparable across methods.

\begin{table}
\centering
\caption{Training performance.}\label{tab:training}
\begin{tabular}{|l|c|}
\hline
Nr of splats & Training time [min]\\
\hline
200k & 30 \\
500k & X \\
1M & Y \\
\hline
\end{tabular}
\end{table}


%\subsection{Ablation Study}
%\textbf{NECESSARY?}


% We analyze the impact of mask supervision on reconstruction quality.

% Specifically, we compare a variant trained using only volumetric intensity supervision to our full model, which additionally incorporates ROI mask guidance.

% Results show that mask supervision improves reconstruction accuracy within the region of interest and leads to better preservation of anatomical structures.


\section{Discussion}\todo{discuss PSNR and other results!}

Our results demonstrate the advantages of reformulating Gaussian Splatting as a direct volumetric optimization problem. By operating directly in voxel space and supervising against a target volume and region-of-interest mask, the method avoids the need for multi-view image rendering, structure-from-motion preprocessing, or camera-based supervision. This provides a more direct optimization pathway that is inherently aligned with the data domain and is particularly well suited for medical imaging applications, where volumetric ground truth is naturally available.

A key advantage of the proposed formulation is that supervision is fully grounded in the native representation of the data. Instead of optimizing against rendered surrogate images, the model is directly constrained by voxel-wise ground truth. This minimizes potential error sources introduced by camera pose estimation, projection geometry, and view-dependent rendering artifacts. In this sense, the proposed approach can be interpreted as a more direct and potentially more precise optimization strategy, which improves robustness and may yield higher reconstruction fidelity compared to image-based Gaussian Splatting methods.

In addition, the representation explicitly stores physically meaningful intensity values per Gaussian primitive (e.g., Hounsfield units in CT data). This design choice enables a major advantage over conventional Gaussian Splatting pipelines: transfer functions can be modified post hoc without requiring retraining or re-optimization of the underlying representation. This stands in contrast to standard 3DGS, where color and appearance are typically static after training. As a result, clinicians or users can interactively adjust contrast, windowing, and tissue emphasis in a manner similar to classical volume rendering workflows, while still benefiting from the efficiency of a sparse Gaussian representation.

Furthermore, initialization can exploit domain knowledge by seeding Gaussians from ROI masks and optionally incorporating anisotropy and structural priors derived from the input volume. This leads to a more task-aware distribution of model capacity compared to standard 3DGS, where initialization is typically less informed and must be inferred indirectly from image observations. The resulting representation serves not only as a rendering primitive but also as a compact volumetric proxy that can be efficiently stored, rendered, and interactively explored.

At the same time, the method introduces new challenges. While dense ray-marched rendering preserves all details of the original data, the Gaussian representation is an approximation and may smooth fine structures or subtle intensity variations. In addition, volumetric rasterization and optimization can be memory- and compute-intensive, especially at large Gaussian counts. Training quality is sensitive to initialization quality, mask accuracy, and hyperparameter selection. Furthermore, the current implementation is not yet fully optimized, and the learned representations may still exhibit relatively uniform, point-cloud-like distributions in certain cases, indicating room for improved adaptive parameterization and densification strategies.

Overall, volumetric Gaussian Splatting occupies a middle ground between dense volume rendering and view-based Gaussian Splatting. It leverages direct voxel-space supervision to improve efficiency and robustness, while reducing reliance on indirect estimation pipelines that can introduce accumulated errors. At the same time, it enables interactive, transfer-function-driven exploration of volumetric data without retraining, making it particularly suitable for medical visualization scenarios where efficiency, controllability, and flexibility are critical.


\section{Conclusion}

We presented a novel formulation of 3D Gaussian Splatting that operates directly in the volumetric domain. By replacing conventional image-based supervision with voxel-space optimization, our approach eliminates the need for camera models and multi-view reconstruction pipelines, resulting in a representation that is naturally aligned with medical imaging data.

The proposed method introduces a differentiable splat-to-volume rasterization and an intensity-aware Gaussian representation, enabling direct optimization of volumetric data and flexible post hoc visualization through transfer functions. This allows Gaussian Splatting to bridge the gap between efficient point-based rendering and classical volume rendering techniques.

Our experiments demonstrate that the resulting representation provides a compelling trade-off between reconstruction quality and computational efficiency. In particular, it enables real-time, volume-rendering-like visualization on resource-constrained hardware, making it well suited for interactive and mixed reality applications.

Future work could focus on improving reconstruction accuracy, further optimizing training performance, and directly computing and storing lighting-related information per splat, such as ambient occlusion, to enable more realistic rendering and to get closer to sophisticated rendering methods like Path Tracing.

We believe that the proposed volume-centric perspective opens new directions for Gaussian-based scene representations beyond traditional image-based settings.

\begin{credits}
\subsubsection{\ackname} A bold run-in heading in small font size at the end of the paper is
used for general acknowledgments, for example: This study was funded
by X (grant number Y).

% \subsubsection{\discintname}
% It is now necessary to declare any competing interests or to specifically
% state that the authors have no competing interests. Please place the
% statement with a bold run-in heading in small font size beneath the
% (optional) acknowledgments\footnote{If EquinOCS, our proceedings submission
% system, is used, then the disclaimer can be provided directly in the system.},
% for example: The authors have no competing interests to declare that are
% relevant to the content of this article. Or: Author A has received research
% grants from Company W. Author B has received a speaker honorarium from
% Company X and owns stock in Company Y. Author C is a member of committee Z.
\end{credits}



\bibliographystyle{splncs04}
\bibliography{bibl}

\end{document}
